\documentclass[journal]{IEEEtran}
\usepackage{amsmath,booktabs,url}
\usepackage[colorlinks=true,urlcolor=blue]{hyperref}
\begin{document}
\title{Supplementary Evidence for Fed-MDBSCAN-G: Audit-v2 Revision}
\author{Working companion to the main manuscript}
\maketitle
\renewcommand{\thetable}{S\arabic{table}}
% FLTrust metric semantics are explained independently of full-method group tables.
\section{Scope and Provenance}
This supplement describes the canonical audit-v2 matrix and diagnostic repetitions separately. Supplementary tables use an independent S-prefixed numbering sequence. Numerical data are copied without changing the canonical run files; SHA-256 hashes are listed in \texttt{evidence/provenance.json}. The scenario manifest enumerates datasets, methods, seeds, attack parameters, and overrides. A reported standard deviation in the CSVs is over seeds, not independent rounds.

The canonical plan contains 2,130 units, 2,125 valid and five failed. A unit denotes a method, condition, and seed. Missing successful JSON output is not evidence that an experiment was never attempted: outcomes distinguish failure from missing work. Table~\ref{tab:failed} lists the persistent failed units. Isolated numerical failures are retained, with no zero-accuracy imputation or seed replacement.
\begin{table}[!t]
\centering\footnotesize
\caption{Canonical failed units. HAR: sample-weighted mean; CIFAR: HDBSCAN FLAME.}\label{tab:failed}
\begin{tabular}{llr}\toprule Dataset & Condition & Seed\\\midrule
HAR & 3.1 & 42\\ HAR & 3.1 & 2024\\ HAR & 3.2 & 137\\ HAR & 3.2 & 2024\\ CIFAR-10 & 7.1 & 42\\\bottomrule
\end{tabular}\end{table}
\section{Temporal Alarm Behavior}
The temporal block activates the attack in zero-based rounds 10--19. Ten preceding and ten following rounds are clean. The table reports final accuracy and post-attack alarm FPR, averaged over three seeds. Optimizer momentum remains unchanged when alarm memory is removed. Persistent alarms after attack removal are false positives, not successful detection of a continuing attack. Identical final accuracies do not imply identical alert behavior.
\begin{table}[!t]
\centering\scriptsize
\caption{Temporal block: accuracy and post-attack alarm FPR (percent).}
\begin{tabular}{lrrr}
\toprule
Data/condition & Variant & Acc. & Post FP \\
\midrule
MNIST 2.1 & Full & 94.07 & 30.0 \\
MNIST 2.1 & No memory & 94.07 & 0.0 \\
MNIST 3.1 & Full & 48.25 & 100.0 \\
MNIST 3.1 & No memory & 48.25 & 100.0 \\
Fashion 2.1 & Full & 83.38 & 30.0 \\
Fashion 2.1 & No memory & 83.38 & 0.0 \\
Fashion 3.1 & Full & 35.44 & 100.0 \\
Fashion 3.1 & No memory & 35.44 & 100.0 \\
HAR 2.1 & Full & 56.43 & 100.0 \\
HAR 2.1 & No memory & 56.81 & 93.3 \\
HAR 3.1 & Full & 41.32 & 100.0 \\
HAR 3.1 & No memory & 41.32 & 96.7 \\
\bottomrule
\end{tabular}
\end{table}
\section{Equal Local-Step Control}
The fixed-step block limits local optimization to five steps. This changes the training protocol and must not be pooled with the epoch-based main matrix. The three compared methods share the same condition and seed within the block. Results below show that equalizing local steps does not uniformly eliminate the utility gap relative to uniform averaging.
\begin{table}[!t]
\centering\scriptsize
\caption{Five-step final accuracy (percent), three seeds per cell.}
\begin{tabular}{lrrr}
\toprule
Data/condition & Full & L0 & Uniform \\
\midrule
MNIST 3.3 & 52.22 & 52.22 & 52.22 \\
MNIST 5.2 & 28.44 & 33.19 & 72.27 \\
Fashion 3.3 & 52.33 & 52.33 & 52.33 \\
Fashion 5.2 & 47.58 & 47.63 & 64.74 \\
HAR 3.3 & 27.17 & 27.17 & 27.17 \\
HAR 5.2 & 44.27 & 44.27 & 44.35 \\
\bottomrule
\end{tabular}
\end{table}
\section{Honest-Client Size Groups}
Client-group false rejection is pooled over observed honest submissions, with the denominator printed explicitly. The size quartile is a recorded partition descriptor, not a demographic category. Counts reuse the same clients across rounds and are not independent observations. They characterize retention disparities; they do not establish causal fairness or minority learning benefit. Dominant-label breakdowns are available in the accompanying CSV.
\begin{table}[!t]
\centering\scriptsize
\caption{Full method, attack-free condition 6.2: size-quartile honest rejection. Q0 is the recorded lowest quartile.}
\begin{tabular}{lrrr}
\toprule
Group & FP & FP+TN & FPR (\%) \\
\midrule
MNIST Q0 & 90 & 2002 & 4.50 \\
MNIST Q1 & 55 & 2047 & 2.69 \\
MNIST Q2 & 1284 & 2027 & 63.34 \\
MNIST Q3 & 2019 & 2024 & 99.75 \\
Fashion Q0 & 45 & 2006 & 2.24 \\
Fashion Q1 & 93 & 2051 & 4.53 \\
Fashion Q2 & 1055 & 2015 & 52.36 \\
Fashion Q3 & 1941 & 2028 & 95.71 \\
HAR Q0 & 10 & 2037 & 0.49 \\
HAR Q1 & 197 & 2019 & 9.76 \\
HAR Q2 & 287 & 2027 & 14.16 \\
HAR Q3 & 1480 & 2017 & 73.38 \\
CIFAR Q0 & 100 & 2007 & 4.98 \\
CIFAR Q1 & 77 & 2044 & 3.77 \\
CIFAR Q2 & 1033 & 2024 & 51.04 \\
CIFAR Q3 & 2025 & 2025 & 100.00 \\
\bottomrule
\end{tabular}
\end{table}
\section{Backend Reproducibility Study}
Normal-profile short repetitions share initial-model and partition hashes but differ at the first client update; 268 of 273 trace entries differ between the two normal repeats. Under deterministic execution, two short repeats and an instrumented short repeat have identical traces. Two further independent 30-round deterministic-profile runs match at all 2,730 trace entries and each ends at accuracy 0.0994.

These runs concern only CIFAR-10, condition 7.1, HDBSCAN FLAME, seed 42. The deterministic profile sets \texttt{CUBLAS\_WORKSPACE\_CONFIG=:4096:8}, enables deterministic PyTorch algorithms, enables cuDNN determinism, and disables cuDNN benchmarking. This is an execution-profile intervention, not a change promoted into the canonical audit-v2 matrix. Trace equality does not prove cross-device repeatability or effective defense. The raw trace and comparison files remain in the investigation directory identified by the project report.


\section{Separate Mechanism Studies}
The following evidence is additional to the canonical 2,130-unit plan. It does not replace failed canonical runs or add independent seeds. A subsequent 117-unit five-step campaign includes 63 attack-free runs; their 1,890 rounds contain 170,100 client submissions. Spearman associations were computed within dataset, method, seed, and round, then summarized by the median over rounds and the median over seeds. Undefined correlations were left undefined. In MNIST, the full method's sample-count/update-norm association is 0.678 overall but 0.056 when restricted to clients with more than 20 samples. Restricting that range changes the population; it does not isolate a causal repair effect. Step equalization does not implement aggregation normalization.

A 27-run single-round pilot and a later shared-checkpoint study use the A/B/C regimes defined in the main text. In the pilot, eight of nine dataset--seed combinations reject no clients in any arm; MNIST/137 is the exception. Therefore a zero median across seeds cannot be interpreted as no rejection in every initial round. The forward study consists of nine reference trajectories and 108 branches, separate from the earlier 117-unit five-step campaign.

\subsection{Matching and Paired Contrasts}
Checkpoint rounds 0, 9, 19, and 29 were specified before the forward branches ran. Models, participation, partitions, and defense history are shared within each cell. Four round-dependent random streams are remapped to the checkpoint round; replay counters are reset while decision memory is restored. Deterministic PyTorch/cuDNN execution and the cuBLAS workspace configuration are used. B/C first-batch and first-update hashes match in all 3,240 client--checkpoint comparisons. All 36 A/reference comparisons match the recorded metrics and norms; reference full-update hashes were not recorded for all 36 cells.

\begin{table}[!t]
\centering\footnotesize
\caption{Selected MNIST checkpoint contrasts. Counts refer to 90 honest participants; FPR is percent. These two cells were selected after inspecting the complete branch study.}\label{tab:branch_cases}
\begin{tabular}{llrrr}\toprule
Seed/round & Arm & L0 rejects & Extra rejects & FPR\\\midrule
137/0 & A & 21 & 0 & 23.33 \\
137/0 & B & 14 & 0 & 15.56 \\
137/0 & C & 5 & 0 & 5.56 \\
2024/9 & A & 20 & 14 & 37.78 \\
2024/9 & B & 15 & 10 & 27.78 \\
2024/9 & C & 14 & 4 & 20.00 \\
\bottomrule\end{tabular}\end{table}

Across the 12 MNIST seed--checkpoint cells, C minus A changes FPR in four cells, with differences from $-17.78$ to $+1.11$ percentage points. The full paired table is supplied in the mechanism evidence directory. These correlated checkpoints and three seeds do not establish an equivalence test or a general gain. Table~\ref{tab:branch_cases} illustrates two selected contrasts, not a representative estimate of deployment benefit. Accuracy at MNIST/2024/round9 changes from 34.51\% to 46.81\% under A to C; this is one-round utility, not a long-horizon learning result.

\subsection{Matched Cluster Replay}
For MNIST/2024/round9, three additional instrumented branches reproduce the original A/B/C input matrices by SHA-256 and reproduce their final decisions. The captured matrices are then passed through the frozen filter offline. Cluster-consensus decisions and SNNC clusters match the instrumented records. Final acceptance can thus be checked directly against the initial accepted pool minus rejected-cluster membership.

A produces six natural clusters, five rejected; only three of those rejected clusters add clients beyond L0, totaling 14. B produces three natural clusters, two rejected, adding ten rejections. C produces two natural clusters, one rejected, adding four. The ten clients avoiding additional rejection under A to C all remain low-density. Four members of one A cluster and four of another become singleton groups; two members of a six-client A cluster also become singletons while the other four remain a rejected cluster.

The neighborhood distance cutoffs, $3\varepsilon$, are 0.208397, 0.178391, and 0.190852 for A/B/C. Natural-cluster size thresholds are 3.230769, 2.073171, and 2, respectively. In C, six of the ten clients have no retained neighbors; the other four form two pairs whose members retain only one another. Mutual neighbors need not share a third neighbor: the two one-element neighbor sets have empty intersection. All ten therefore remain singleton groups, below the cluster-size threshold, and are not subject to additional cluster rejection.

The three A clusters responsible for additional rejection have centroid-distance/consensus-threshold ratios 1.837, 1.895, and 1.990. C's remaining rejected cluster has ratio 1.718. Hence the observed change is not those ten clients joining a cluster that now passes consensus. All arms keep the alarm active and use the L0+L2 path without the safety valve. The six other newly accepted clients change their L0 status. The two layer counts are descriptive accounting, not independent causal mediation effects.

\subsection{Comparison with the Original MDBSCAN Specification}
Inspection of the original article (page 5, Algorithms 1--2; page 6, Fig.~7) clarifies three differences. Its SNNC pseudocode does not specify the additional $3\varepsilon$ neighbor cutoff; its natural-cluster threshold is based on mean group size without the explicit floor of two used locally; and its remaining points undergo DBSCAN followed by nearest-cluster assignment for DBSCAN noise. The federated filter does not execute that last stage. A separate generic clustering routine exists in the source, but it is not the server's defense path. The filter's \texttt{min\_pts} argument does not enter its acceptance computation.

There is also a boundary ambiguity in the article itself: Eq.~(4) uses a strict greater-than comparison against mean group size, whereas Algorithm 1, line 29, and Fig.~7 use greater-than-or-equal. A reproduction should state the chosen interpretation and check equality cases. The local implementation uses greater-than-or-equal with its added floor. This textual comparison is not verification against the original authors' executable implementation. Consequently, the observed singleton acceptance is attributed to the local federated adaptation, not to the original clustering algorithm.

\subsection{Interpretation Limits}
Batching changes updates, neighbor rankings, radius estimates, and the initial trusted pool together. No intervention here isolates radius from update geometry. Likewise, whole-client arrival loss and post-training loss are not measured by the first and fifth training-mini-batch losses: the fifth loss precedes the fifth optimizer step, and its batch may differ from the first. Net update norm is distinct from cumulative path length and from distance to the cohort's geometric median.

The diagnosis concerns the local SNNC adaptation and a selected attack-free checkpoint. It neither verifies fidelity to an original author's implementation nor demonstrates an adversarial exploit. Repeating one batch is not proposed as a defense. Candidate corrections require designated development conditions and independent validation of clean retention, utility, and attack success.

\section{Fixed-Geometry Unclustered-Update Policies}
We retain the original neighborhood cutoff and evaluate three acceptance policies at the 72 development matrices under each of the original and density-only gates: $72\times2\times3=432$ offline evaluations. Three historical A/B/C matrices add nine evaluations, reported separately. All 147 baseline branches reproduce their recorded accepted IDs and filter information exactly. The 441 evaluations are not training runs or independent replicates.

Let $B_0$ be the L0 accepted set, $L$ the executed low-density partition, and $C$ the union of natural-cluster members. The intervention targets $E=B_0\cap(L\setminus C)$. P0 retains the original policy. P1 applies the frozen cluster-consensus function to each member of $E$ as a singleton, using the unchanged $B_0$ and threshold factor two; failing singletons are removed in addition to the original cluster rejections. P2 removes all of $E$ without a consensus test. Each policy then applies the original safety valve: if fewer than half the participating updates remain, it restores $B_0$. If L2 does not execute, all policies retain the original L0 decision. Geometry is held fixed within each comparison, and labels are used only for reporting.

\begin{table}[!t]
\centering\scriptsize
\caption{Fixed-geometry policy outcomes. Each triple is P0/P1/P2. FP and TP count final honest and attacker rejections; V counts valve activations. Each row covers 18 checkpoints. Attack rows contain 1296 honest and 324 attacker observations; clean rows contain 1620 honest observations and no attackers. Gates reuse the same matrices.}\label{tab:unclustered_policy}
\setlength{\tabcolsep}{3pt}
\begin{tabular}{llrrr}\toprule
Gate & Condition & FP & TP & V\\\midrule
Original & Min-Max attack & 0/0/0 & 0/0/0 & 0/0/0 \\
Original & Min-Max clean & 319/351/326 & --- & 1/5/8 \\
Original & Patch attack & 0/0/0 & 0/0/0 & 0/0/0 \\
Original & Patch clean & 1/1/1 & --- & 0/0/0 \\
Density & Min-Max attack & 0/0/0 & 0/0/0 & 0/0/18 \\
Density & Min-Max clean & 319/395/404 & --- & 1/5/9 \\
Density & Patch attack & 0/1/142 & 0/0/46 & 0/0/0 \\
Density & Patch clean & 1/1/82 & --- & 0/0/0 \\
\bottomrule\end{tabular}
\end{table}

Under the original gate, L2 executes in 10 of 72 checkpoints and in none of the 36 attacked checkpoints. P1 changes final accepted identities in four checkpoints, adding 32 honest rejections. Under the density-only gate, L2 executes in 66 checkpoints; P1 changes eight, adding 76 honest rejections in clean conditions and one in an attacked patch checkpoint. Neither gate yields an additional attacker rejection under P1.

P2 illustrates the valve interaction. In the density-only Min-Max attack block it removes 1296 honest observations before the valve; the attackers are outside the targeted low-density set. The valve restores every update in all 18 checkpoints. In the density-only patch attack block, P2 instead rejects 46 attacker and 142 honest observations without triggering the valve. These counts do not measure attack success or learning utility. In the clean Min-Max block, P2's net honest-rejection increases of seven and 85 under the original and density-only gates comprise 11 and 89 newly rejected observations, respectively, offset by four newly accepted observations in each gate. A stricter pre-valve rule need not produce a subset of the baseline final accepted set.

For the historical MNIST/2024/round9 A/B/C case, $|E|=12/23/31$. P1 adds 10/19/26 honest rejections, changing final counts from 34/25/18 to 44/44/44 without triggering the valve. Equal counts do not imply identical identities: A and B have the same P1 accepted set, while C differs by one identity swap. All ten previously identified clients escaping A's additional rejection in C fail the C singleton test. P2 leaves only 44/42/41 accepted updates before the valve, below the threshold of 45; it therefore restores $B_0$, leaving 20/15/14 final rejections.

Input and frozen-source hashes are checked before and after analysis. Each frozen module passes seven boundary checks, including threshold equality and the valve boundary. The original consensus function is called with an exactly checked cache of its geometric median; the first eligible singleton in each branch is also compared with an uncached call. A separate accounting check verifies all 441 rows against their saved identity sets. The inputs were already used in development; the protocol hash establishes content identity, not independent preregistration. These interventions diagnose the local acceptance path and do not establish a corrected defense.

\section{Evidence File Guide}
\begin{itemize}
\item \texttt{units.csv}: successful-unit final metrics and identity fields.
\item \texttt{outcomes.csv}: all planned unit outcomes, including failures.
\item \texttt{cells.csv}, \texttt{cell\_status.csv}: cell summaries and coverage; incomplete cells must retain their success counts.
\item \texttt{paired\_deltas.csv}, \texttt{paired\_summary.csv}: identity-matched method/control contrasts; contrast direction is explicit.
\item \texttt{benign\_groups.csv}: honest-client rejection counts by group.
\item \texttt{scenario\_manifest.json}: exact 149-job plan.
\item \texttt{provenance.json}: source identity and copied CSV hashes.
\end{itemize}
The separate \texttt{evidence/mechanism\_20260914/} directory contains the stratified summary, all paired checkpoint contrasts, layer counts, selected cluster membership, neighborhood diagnostics, and their provenance. Canonical provenance remains unchanged. The separate \texttt{analysis/unclustered\_policy\_20260920/} directory supplies per-cell and grouped policy counts, the execution protocol, analysis code, and input/output hashes. Identity-level outputs require the local raw archive.

The preserve-empty and oracle blocks remain in these files. They are protocol sensitivity and evaluator-reference experiments, respectively; neither is silently pooled into the main method ranking. Full statistical inference, original-author baseline reproduction, public artifact deposition, and a complete novelty assessment remain author-review tasks before submission.
\section{FLTrust Metric Interpretation and Fidelity Scope}
For the normalized FLTrust implementation, stored client FPR and TPR count zero-trust-weight assignments to honest and malicious updates in the ordinary cosine-weighting path. They are not evidence of a separate malicious-client classifier. The full-method group-rejection tables in this supplement retain their original meaning. Missing or near-zero root updates invoke a local fallback, so the zero-weight interpretation must not be extended to fallback rounds without checking the executed path.

The formula comparison uses 18 recorded matrices, four constructed root vectors, and two local variants: 144 comparison rows, with 72 input cases for each variant. It compares a NumPy transcription of the author-demo rule with a helper using the local filter and a re-expression of server aggregation. That initial comparison does not execute the author training system or the production server path. A subsequent check passes a subset of six matrices with the four root constructions through the actual local server, with relative aggregate discrepancy below $2.4\times10^{-7}$. Boundary probes also show divergent fallback behavior: a degenerate root produces an unflagged uniform mean, while all-zero trust invokes the configured degraded accept-all fallback. In the canonical archive, 135 of 5490 rounds accept everyone; missing root norms and executed-operator records prevent exclusion of the silent root fallback in those rounds. These checks do not establish equivalence of root and client training schedules, whose optimizer-step counts can differ. Only the normalized variant is identified as canonical in the accompanying audit; the clipped-only variant should not be interpreted as the published normalization rule.
\section{Scope of the FLAME Path Checks}
Six synthetic input cases exercise the production server aggregation: a majority cluster, a nonzero global model, identical updates, two clients, zero updates, and a forced empty clustering result. The last case replaces only clustering labels; it is a control-flow check, not an observed frequency. A deterministic noise spy records the Gaussian call's parameters and returns a known vector. Conditional on the selected IDs, independently computed clipping, averaging, and noise agree within $6.65\times10^{-8}$ maximum absolute error. This checks arithmetic and noise placement, not Gaussian sampling or author-code clustering equivalence. The local FLAME filter and aggregation-branch ASTs match their frozen canonical counterparts. The accompanying audit retains hashes and per-run coverage for 200 complete runs and one failed run.
\end{document}
